\documentclass[11pt]{article}

\usepackage[T1]{fontenc}
\usepackage{amsmath,amssymb}
\usepackage{hyperref}

\title{Reduced-Bias Estimation for ML, FIML, and Moment-Quadratic SEM}
\author{}
\date{\today}

\newcommand{\T}{^{\mathsf T}}
\newcommand{\tr}{\operatorname{tr}}
\newcommand{\E}{\mathbb{E}}
\newcommand{\hatth}{\widehat\theta}
\newcommand{\hatalpha}{\widehat\alpha}
\newcommand{\vech}{\operatorname{vech}}

\begin{document}
\maketitle

\section*{Purpose}

This note records the common estimating-equation convention used by the
\texttt{estimate::frontier::rbm\_*} implementation. The first implementation
covered complete-data ML and FIML. The extension adds continuous ULS/GLS/WLS,
all-ordinal LS, mixed ordinal LS, and ML2S. All paths use the same
linear-constraint-reduced trace algebra; they differ only in how they supply
the observed bread and the empirical estimating-function meat.

\section*{Objective scale}

Let the estimator minimize a per-observation objective
\[
  f_N(\theta) = \frac{1}{N}\sum_{i=1}^N f_i(\theta),
  \qquad
  U_N(\theta) = \frac{\partial f_N(\theta)}{\partial \theta}.
\]
For likelihood estimators \(f_i\) is the normal deviance contribution up to the
usual SEM scale. For moment-quadratic estimators
\[
  f_N(\theta)
    =
  \frac12 \sum_b \frac{n_b}{N}
    d_b(\theta)\T W_b d_b(\theta),
  \qquad
  d_b(\theta) = m_b(\theta) - s_b .
\]
The implementation stores the observed bread on the total-\(N\) scale,
\[
  J_N(\theta) =
    N\,\frac{\partial U_N(\theta)}{\partial\theta\T},
\]
and the meat as the empirical sum
\[
  E_N(\theta) = \sum_{i=1}^N \psi_i(\theta)\psi_i(\theta)\T,
  \qquad
  \sum_i \psi_i(\hatth) \approx 0 .
\]
With this convention the trace term is dimensionless:
\[
  T(\theta) = \tr\{J_N(\theta)^{-1}E_N(\theta)\}.
\]

\section*{Linear constraints}

The lavaanified model may contain linear equality constraints. Write the
free-parameter vector as
\[
  \theta = \theta_0 + K\alpha,
\]
where the columns of \(K\) span the feasible tangent space. RBM is carried out
in \(\alpha\)-space:
\[
  J_\alpha = K\T J_N K,
  \qquad
  E_\alpha = K\T E_N K,
  \qquad
  T_\alpha = \tr(J_\alpha^{-1}E_\alpha).
\]
The explicit one-step correction computes the gradient of
\[
  P(\alpha) = -\frac12 T_\alpha(\alpha)
\]
at the unadjusted solution and applies
\[
  \Delta_\alpha = J_\alpha^{-1}
    \frac{\partial P(\alpha)}{\partial\alpha},
  \qquad
  \hat\theta_{\mathrm{eRBM}} =
    \hat\theta + K\Delta_\alpha .
\]
The implicit estimator minimizes the adjusted objective
\[
  f_N(\theta) - \frac{1}{N}P(\theta)
    =
  f_N(\theta) + \frac{1}{2N}T_\alpha(\theta).
\]
This sign convention matches the metadata stored by the implementation:
\texttt{penalty = -0.5 * trace\_term} and
\texttt{penalized\_fmin = fmin + penalty\_per\_observation}.

\section*{ML and FIML}

For complete-data ML, \(\psi_i\) is the row-wise gradient of the normal
log-likelihood contribution in the full free-parameter basis. For FIML,
\(\psi_i\) is the observed-pattern score contribution for the row's observed
subvector. The FIML complete-data special case is therefore pinned to the ML
path by a regression test: after scale normalization, bread, meat, trace, and
the explicit correction agree.

\section*{Moment-quadratic estimators}

For a moment block, let
\[
  D_b(\theta) = \frac{\partial m_b(\theta)}{\partial\theta\T}.
\]
The stage-two estimating equation has the form
\[
  U_N(\theta) =
    \sum_b \frac{n_b}{N} D_b(\theta)\T W_b d_b(\theta).
\]
The observed bread is the derivative of this estimating equation, not only the
Gauss-Newton term. In the continuous LS path this is the same observed-Hessian
bread used by the misspecification-robust SE implementation. Ordinal, mixed,
and ML2S paths reuse the corresponding weighted-inference observed bread.

The meat is built from complete infinitesimal-jackknife rows for the fitted
moment estimator. In block notation the implementation computes rows whose
crossproduct is equivalent to
\[
  E_\alpha =
    \sum_b (D_b K)\T V_b\T V_b (D_b K),
\]
where \(V_b\) is the casewise transformed moment influence, including any
estimated-weight or first-stage contribution. This is why ordinal DWLS/WLS,
mixed polyserial/polychoric LS, DLS/ADF, and ML2S do not use a stage-two-only
meat: their first-stage saturated-moment and weight estimation effects are
already in \(V_b\).

\section*{Estimator routing}

The common RBM result reports both full-\(\theta\) and reduced-\(\alpha\)
bread/meat diagnostics. Current routed estimators are:
\[
\begin{array}{ll}
\text{ML} & \text{complete-data normal scores},\\
\text{FIML} & \text{observed-pattern normal scores},\\
\text{ULS/GLS/WLS} & \text{continuous moment IJ rows},\\
\text{ordinal ULS/DWLS/WLS} & \text{polychoric/threshold IJ rows},\\
\text{mixed ordinal ULS/DWLS/WLS} & \text{mixed moment IJ rows},\\
\text{ML2S} & \text{FIML saturated-moment IJ rows with NT/DWLS/ADF/DLS weights}.
\end{array}
\]
Nonlinear equality constraints remain outside the implementation because the
current correction requires a fixed linear tangent matrix \(K\) over the local
finite-difference stencil.

\end{document}
